% ICLR 2027 submission (analysis paper).
\documentclass{article}
\usepackage{iclr2027_conference,times}
\input{math_commands.tex}
\usepackage{hyperref}
\usepackage{url}
\usepackage{amsmath,amssymb,amsfonts}
\usepackage{graphicx}
\usepackage{booktabs}
\usepackage{microtype}
\usepackage{multirow}
\usepackage{xcolor}
\usepackage{algorithm}
\usepackage{algpseudocode}

\newcommand{\model}{\textsc{Qwen3}}
\newcommand{\probe}{\textsc{LinearProbe}}
\newcommand{\feat}{\textsc{Feat}}
\newcommand{\llmprobe}{\textsc{FrozenLLM}}
\newcommand{\tsp}{\textsc{TS-Encoder}}

\title{Recognize, Don't Generate:\\ Language Pretraining as a Temporal-Structure
Prior for Time-Series Forecasting}

\author{Anonymous Authors}

\date{}

\begin{document}
\maketitle

\begin{abstract}
Whether language pretraining helps time-series forecasting is disputed: some
studies report large gains while others find that pretrained language models
offer no benefit over random initialization or specialized architectures.  We
argue that both sides are partially right because the field has conflated two
distinct capabilities: \emph{recognizing} temporal structure and
\emph{generating} future numeric values.  Using a tightly controlled protocol
in which the only variable is the model initialization, we first confirm that
language pretraining substantially improves forecasting under LoRA
fine-tuning.  We then isolate the source of this benefit with frozen-model
probes: a frozen pretrained LLM recognizes trend, periodicity, local
autoregressive, and regime-switching dynamics far more accurately than a random
same-architecture control, and---critically---this recognition transfers
out-of-distribution to dynamics never seen during probe training, where
hand-crafted features, supervised patch transformers, and MLPs all collapse.
On the very same windows, however, frozen greedy numeric generation is far less
reliable: low parse coverage and several-fold higher error than the oracle
expert, a gap that persists even after grammar-constrained decoding removes
formatting failure.  This recognize-don't-generate asymmetry holds across model
scales (0.6B--8B) and yields a practical mechanism: a frozen LLM used purely as
a temporal-expert router improves zero-shot routing on 15 real-world datasets
and enables patch-level dynamic routing.  Our results reframe the role of
language pretraining in time-series modeling: it is a reusable
\emph{temporal-structure prior}, best exploited through recognition and routing
rather than direct numeric generation.
\end{abstract}

\section{Introduction}\label{sec:intro}

Large language models (LLMs) are increasingly applied to time-series
forecasting, either as backbones fine-tuned on numerical text
\citep{gruver2023large,time-llm2023,qwen-time} or as frozen feature
encoders \citep{liu2024pre-training,rasul2024lag}.  The central empirical
question---\emph{does language pretraining actually help forecasting?}---has
produced sharply conflicting answers.  \citet{tan2024are} report that
pretrained LLMs do not consistently beat scratch models under their protocol,
while several 2025--2026 studies \citep{randominit2025,manifold2026} find that
language pretraining provides a measurable, sometimes large, advantage.  These
discrepancies are typically attributed to protocol differences (tokenizer,
backbone, fine-tuning strategy, data leakage), making the literature difficult
to interpret.

We argue that the conflict is more fundamental.  Forecasting with an LLM
requires two very different abilities: (i) \emph{recognizing} the temporal
structure of the observed history (trend, periodicity, volatility, regime) and
(ii) \emph{generating} numeric values that continue that structure.  Most
prior work treats these as one capability and evaluates the end-to-end
forecast, so a failure of either masks the other.  In this paper we
disentangle them under a tightly controlled protocol and find a striking
asymmetry:

\begin{center}
\fbox{\parbox{0.92\linewidth}{\centering
Language pretraining transfers temporal \emph{recognition} robustly, even
out-of-distribution, but does \emph{not} confer reliable numeric
\emph{generation}.}}
\end{center}

Concretely, we make the following contributions.
\begin{itemize}
  \item \textbf{Controlled causal evidence for pretraining benefit (E1).}
  Holding architecture, tokenizer, windowing, and fine-tuning protocol
  identical, and varying only the initialization (pretrained vs.\ random), we
  confirm that language pretraining improves LoRA-finetuned forecast MSE by
  $2$--$25\times$ on synthetic and real data, and that this is not a LoRA
  artifact: with full fine-tuning on small Qwen3 models, random initialization
  still lags $10$--$13\times$ behind.
  \item \textbf{Frozen recognition is robust and transfers OOD (E2, E4).}
  A frozen pretrained LLM linearly separates five temporal dynamics with
  $\sim$96--98\% accuracy vs.\ $\sim$79--91\% for a random same-architecture
  control.  When the probe is trained only on three clean dynamics and tested
  on unseen mixtures and regime switches, the pretrained representation
  transfers ($67.8\%$ oracle-hit) while hand features ($12.8\%$), random
  controls ($28.9\%$), a supervised patch transformer ($13.3\%$), and an MLP
  ($34.2\%$) all collapse.
  \item \textbf{Recognition $\gg$ generation on identical windows (E3).}
  On the same windows, frozen recognition reaches $96$--$98\%$ while frozen
  greedy generation has parse rates of $27$--$98\%$ and $5$--$10\times$ higher
  MSE than the oracle expert.  Grammar-constrained decoding raises parse rate
  to $100\%$ but the error gap persists ($\approx 6\times$), ruling out a mere
  formatting artifact.
  \item \textbf{Scale and realism (R5, R6).}  The asymmetry holds from
  $0.6$B to $8$B parameters, and the recognition-based router improves
  zero-shot expert selection on 15 real-world datasets (ETT, Electricity,
  Weather, Exchange, Traffic, Illness, M4) and enables patch-level dynamic
  routing that improves regime-switching windows by $33\%$ while honestly
  remaining neutral on homogeneous windows.
\end{itemize}

\section{Related Work}\label{sec:related}

\textbf{LLMs for time-series forecasting.}  A large body of work fine-tunes or
prompts LLMs for forecasting \citep{gruver2023large,time-llm2023,qwen-time,
new2022,zhou2024one}.  \citet{tan2024are} provide a controlled re-analysis and
find that pretrained LLMs do not consistently help, questioning the language
prior.  In contrast, \citet{randominit2025} show that with matched protocols,
language transfer yields a persistent advantage, and \citet{manifold2026}
analyze Qwen3 and argue that pretraining creates a reusable sequential
manifold.  Our paper sits between these: we confirm the advantage but show
that its mechanism is \emph{recognition}, not \emph{generation}.

\textbf{Probing and analysis of pretrained models.}  Linear probing is a
standard tool for studying representations \citep{alain2016understanding,
beldon2018measuring}.  Recent work probes LLMs on temporal data
\citep{liu2024pre-training,tan2024are}.  \citet{manifold2026} probe Qwen3 on
synthetic dynamics; our layer-wise probes (Section~\ref{sec:layerwise})
complement this by showing where in the stack the temporal structure lives and
how forecasting fine-tuning shifts it.  Parallel work on time-series
foundation models probes \emph{temporal concepts} (trend, periodicity, level)
inside their representations and studies representation $\to$ concept $\to$
intervention chains \citep{tsfm_probe2025,tsfm_intervention2025}.  We share
the representation-probing methodology but ask a different attribution
question: instead of asking \emph{which} temporal concepts are encoded, we ask
\emph{whether the language-pretraining prior transfers through the
recognition interface but not the native numeric-generation interface}, and
whether the accessible structure is compositionally useful for routing.  Our
key observation---that the \emph{same} frozen backbone gives opposite
conclusions under different evaluation interfaces (linear probe vs.\ native
decoding)---is the contrast that distinguishes this work from prior probing
analyses.

\textbf{Mixture-of-experts and routing for time series.}  Routing among
specialized forecasters dates back to forecast combination
\citep{bates1969combination}.  Modern MoE \citep{shazeer2017outrageously}
routes tokens to experts; here we route \emph{temporal dynamics} to
\emph{temporal experts}, using the LLM as a structure recognizer rather than a
forecaster, echoing the ``Task$\to$Router$\to$Task'' finding of
\citet{waveTLM2025}.

\section{Controlled Protocol}\label{sec:protocol}

\paragraph{Dynamics.}  We generate labeled windows from five dynamics:
\emph{trend} (linear slope $+$ noise), \emph{periodic} (sine with period
$12$ or $24$), \emph{local} (stable AR(1) with $\phi\in[0.7,0.95]$),
\emph{mixture} (trend $+$ sine), and \emph{regime} (trend in the first $60\%$
of the context, sine afterward).  Each window has a 64-step context and a
16-step future, standardized by the context.  The first three are ``clean''
and the last two are reserved for out-of-distribution (OOD) testing.

\paragraph{Text protocol.}  A frozen LLM reads the context as
``\texttt{History: +0.12 -0.04 ... Forecast:}'' with values clamped to
$[-9.99,9.99]$ and two decimals.  The final-token hidden state is the
representation.

\paragraph{Probe.}  A linear softmax probe is trained on train-window
representations and evaluated on held-out test windows.  The feature floor uses
nine hand-crafted statistics (slope, linear $R^2$, volatility, lag-1
autocorrelation, dominant period/strength, spectral entropy, spectral peak
concentration, range ratio).

\paragraph{Experts and router.}  Three lightweight univariate experts:
\emph{trend} (robust linear extrapolation), \emph{periodic} (autocorrelation
period detection $+$ seasonal naive), and \emph{local} (damped AR(5)).  For
every window the \emph{oracle} expert is the one with the lowest forecast MSE.
A router predicts the expert from the representation (probe) or from features.

\paragraph{Models.}  Qwen3-0.6B / 1.7B (vanilla) and the Wave-MoE Qwen3-8B
backbone (language trunk only).  Random controls use the identical
architecture with re-initialized weights.

\section{Results}\label{sec:results}

\subsection{Language pretraining helps forecasting (E1)}
Table~\ref{tab:e1} summarizes the controlled ablation.  With the only variable
being initialization, pretrained checkpoints reduce test MSE by $2.2$--$25\times$
across datasets and seeds, with far lower variance and $100\%$ parse coverage.
The small-model $2\times2$ (size $\times$ LoRA/full-FT) control confirms the
gap is not a low-rank artifact: across three seeds of full fine-tuning on
Qwen3-0.6B, pretrained initialization yields lower MSE in all three runs, but
random full fine-tuning is highly variable (per-seed ratios $23.7\times$,
$1.7\times$, $46.5\times$; ratio of mean MSEs $20.9\times$).  This is not
evidence that pretrained weights already forecast numerically---it shows they
provide a substantially better starting point for learning the forecasting
map, which motivates our separation between accessible structure and native
numeric generation below.

\begin{table}[t]
\centering
\small
\begin{tabular}{lccc}
\toprule
Dataset & Pretrained MSE & Random MSE & Reduction \\
\midrule
Synthetic sine & $0.0094\pm0.0006$ & $0.0212\pm0.0011$ & $2.2\times$ \\
ETTm1 & $0.0099\pm0.0017$ & $0.2453\pm0.1366$ & $25\times$ \\
ETTh1 & $0.0487\pm0.0097$ & $0.2896\pm0.2059$ & $6\times$ \\
\bottomrule
\end{tabular}
\caption{E1: initialization-only ablation, LoRA, 3 seeds.}
\label{tab:e1}
\end{table}

\subsection{Frozen recognition is accurate and scales (E2, R5)}
Table~\ref{tab:scale} shows recognition accuracy, generation reliability, and
zero-shot routing across model sizes.  Three findings stand out:
(i) frozen recognition is high and stable ($\geq96\%$) at every scale, well
above the random controls; (ii) generation parse rate \emph{decreases} with
scale while generation error stays $5$--$10\times$ the oracle expert; (iii)
the routing advantage of pretraining is largest at $8$B ($67.8\%$ vs.\ random
$28.9\%$, features $12.8\%$).

\begin{table}[t]
\centering
\small
\begin{tabular}{llccccc}
\toprule
Model & Family & Recog. (pre.) & Recog. (rand.) & Gen. parse & Gen. / oracle & Router (pre.) \\
\midrule
DistilGPT2 & GPT & $0.973\pm0.009$ & $0.744\pm0.022$ & $1.000$ & $15.4\times$ & $0.181\pm0.037$ \\
GPT-2 (124M) & GPT & $0.980\pm0.008$ & $0.699\pm0.041$ & $1.000$ & $12.4\times$ & $0.194\pm0.063$ \\
Qwen3-0.6B & Qwen & $0.984\pm0.006$ & $0.789\pm0.010$ & $0.980$ & $9.9\times$ & $0.337\pm0.068$ \\
Qwen3-1.7B & Qwen & $0.979\pm0.011$ & $0.837\pm0.012$ & $0.575$ & $8.1\times$ & $0.304\pm0.012$ \\
Qwen3-8B & Qwen & $0.963\pm0.005$ & $0.906\pm0.009$ & $0.267$ & $5.5\times$ & $0.678\pm0.075$ \\
\bottomrule
\end{tabular}
\caption{R5 family $\times$ scale: recognition, generation, routing (3 seeds,
mean$\pm$std).  Feature floor is $0.99$ for all.  GPT-2/DistilGPT2 parse $100\%$
of outputs yet are $12$--$15\times$ worse than the oracle expert---the largest
generation gaps of any family---confirming that generation failure is a
continuation/regression failure, not a formatting one.}
\label{tab:scale}
\end{table}

\paragraph{Cross-family confirmation.}  The recognize-don't-generate pattern
is not specific to Qwen: GPT-2 and DistilGPT2 show the same asymmetry, with an
even \emph{larger} pretraining advantage on recognition ($97$--$98\%$ vs.\
$70$--$74\%$ random) and the largest generation gaps of all families
($12$--$15\times$ oracle despite $100\%$ parse rate).  This is strong evidence
that the effect is a general property of language pretraining, not of a
particular tokenizer or architecture.

\subsection{Layer-wise structure of the temporal prior (R)}
\label{sec:layerwise}
Figure~\ref{fig:layerwise} plots per-layer probe accuracy for frozen
pretrained vs.\ random Qwen3-0.6B.  Two qualitative differences stand out.
First, temporal structure is distributed across the stack in the pretrained
model: recognition rises from chance at the embedding layer to $\sim$93--97\%
by layer 1 and plateaus at $\sim$98\% through the upper layers.  Second, the
random same-architecture control peaks early (layer 2--3, $\sim$90\%) and then
\emph{decays} with depth (down to $\sim$79\% at the final layer).  Language
pretraining therefore does not merely add a single ``temporal head'': it
preserves linearly separable temporal structure throughout the depth of the
network, whereas in a randomly initialized network that structure is confined
to the shallowest layers.

\begin{figure}[t]
\centering
\fbox{\parbox{0.9\linewidth}{\centering\vspace{2.2cm}
\textcolor{gray}{Layer-wise probe accuracy: pretrained (plateau $\approx$0.98)
vs.\ random (peak $\approx$0.90 at layer 2--3, decaying to $\approx$0.79).}}}\vspace{-0.4cm}
\caption{Per-layer linear-probe accuracy on the five dynamics.}
\label{fig:layerwise}
\end{figure}

\subsection{Recognize vs.\ generate on identical windows (E3)}
Table~\ref{tab:e3} shows the paired comparison: the \emph{same} windows are
classified by the frozen probe and forecast by frozen greedy generation.  The
recognition accuracy ($96$--$98\%$) contrasts sharply with generation: parse
coverage as low as $27\%$ (8B) and MSE $5$--$10\times$ the oracle expert even
when parsed.  Grammar-constrained decoding raises parse rate to $100\%$ but
the MSE gap remains $\approx 6\times$, ruling out a formatting-only
explanation (Table~\ref{tab:e3b}).

\begin{table}[t]
\centering
\small
\begin{tabular}{lccccc}
\toprule
Model & Recog. acc. & Gen. parse & Gen. MSE & Oracle MSE & Ratio \\
\midrule
Qwen3-0.6B & $0.98$ & $0.98$ & $3.31$ & $0.34$ & $9.9\times$ \\
Qwen3-1.7B & $1.00$ & $0.58$ & $2.73$ & $0.34$ & $8.1\times$ \\
Qwen3-8B & $0.97$ & $0.27$ & $1.85$ & $0.33$ & $5.5\times$ \\
\bottomrule
\end{tabular}
\caption{E3: paired recognize-vs-generate on identical windows.}
\label{tab:e3}
\end{table}

\begin{table}[t]
\centering
\small
\begin{tabular}{lcc}
\toprule
Decoding & Parse rate & MSE / oracle \\
\midrule
Free-form greedy & $0.267$ & $5.5\times$ \\
Grammar-constrained & $1.000$ & $6.2\times$ \\
\bottomrule
\end{tabular}
\caption{E3b: constrained decoding removes formatting failure; the gap
persists.}
\label{tab:e3b}
\end{table}

\paragraph{The generation gap survives every prompt and decoding choice.}
Table~\ref{tab:genrob} attacks the ``you didn't prompt it properly'' objection
by varying sampling (greedy vs.\ temperature $0.7/1.0$), adding 1--2 in-context
examples, and adding an explicit ``only output numbers'' instruction, all on
the \emph{same} stratified $40$-window evaluation set as E3.  For both the
$0.6$B and $8$B models, generation stays $4.9$--$9.9\times$ worse than the
oracle expert across every variant: sampling does not help (it lowers parse
coverage while the error ratio stays high), few-shot demonstrations raise
parse coverage but leave the error ratio $\geq 4.9\times$, and the explicit
instruction hurts both models' coverage.  Combined with the
grammar-constrained result (parse $=100\%$, ratio $6.2\times$), the
recognition--generation gap is not attributable to formatting, decoding, or
prompting---it reflects a genuine deficit in numeric continuation.

\begin{table}[t]
\centering
\small
\begin{tabular}{lccccc}
\toprule
 & \multicolumn{2}{c}{Qwen3-0.6B} & \multicolumn{2}{c}{Qwen3-8B} \\
\cmidrule(lr){2-3}\cmidrule(lr){4-5}
Decoding / prompt & Parse & MSE/oracle & Parse & MSE/oracle \\
\midrule
Greedy & $1.000$ & $9.6\times$ & $0.275$ & $5.9\times$ \\
Temp $0.7$ & $0.950$ & $7.1\times$ & $0.375$ & $8.0\times$ \\
Temp $1.0$ & $0.925$ & $8.4\times$ & $0.300$ & $9.7\times$ \\
Instruction & $0.925$ & $8.8\times$ & $0.475$ & $8.5\times$ \\
1-shot & $1.000$ & $8.3\times$ & $0.925$ & $4.9\times$ \\
2-shot & $1.000$ & $9.9\times$ & $0.875$ & $7.4\times$ \\
\bottomrule
\end{tabular}
\caption{E3c: generation error vs.\ oracle expert under every prompt/decoding
variant on the same stratified windows as E3.  The $4.9$--$9.9\times$ gap
persists across all of them.}
\label{tab:genrob}
\end{table}

\subsection{The asymmetry persists across task difficulty (R)}
Table~\ref{tab:asym} sweeps context length, noise, and horizon for the frozen
$0.6$B model on identical windows.  Recognition remains $>90\%$ even under
heavy noise, and \emph{improves} with horizon ($96.7\%$ at $H{=}8$ to $99.2\%$
at $H{=}32$).  In contrast, frozen generation stays $6$--$9\times$ worse than
the oracle expert and its relative gap \emph{grows} with horizon---exactly the
error-accumulation signature predicted in Section~\ref{sec:analysis}.

\begin{table}[t]
\centering
\small
\begin{tabular}{lcccc}
\toprule
Axis & Setting & Recog. acc. & Gen. parse & Gen. / oracle \\
\midrule
\multirow{3}{*}{Context} & $32$ & $0.942$ & $0.975$ & $5.8\times$ \\
 & $64$ & $0.983$ & $1.000$ & $7.9\times$ \\
 & $128$ & $0.983$ & $1.000$ & $8.6\times$ \\
\midrule
\multirow{3}{*}{Noise} & $0.05$ & $0.983$ & $1.000$ & $7.9\times$ \\
 & $0.10$ & $0.958$ & $1.000$ & $7.8\times$ \\
 & $0.20$ & $0.908$ & $1.000$ & $7.1\times$ \\
\midrule
\multirow{3}{*}{Horizon} & $8$ & $0.967$ & $1.000$ & $6.5\times$ \\
 & $16$ & $0.983$ & $1.000$ & $7.9\times$ \\
 & $32$ & $0.992$ & $0.875$ & $8.4\times$ \\
\bottomrule
\end{tabular}
\caption{Difficulty sweeps: recognition robust, generation gap grows.}
\label{tab:asym}
\end{table}

\subsection{Zero-shot routing: recognition transfers OOD (E4)}
The probe is trained on the three clean dynamics and evaluated on unseen
mixture/regime windows.  Because the raw mixture/regime set is dominated by
trend-oracle windows ($\sim$87\%; local-oracle windows are structurally
absent), raw oracle-hit accuracy and overall MSE are misleading and
favourable to a trivial constant-trend policy.  We therefore treat this set as
a diagnostic (Table~\ref{tab:e4}) and use a \emph{balanced} OOD test as the
primary evidence (Table~\ref{tab:e4bal}): we rejection-sample mixture/regime
windows so that the trend-oracle and periodic-oracle classes are equally
represented (60 each, per seed, no OOD tuning).

On the balanced set, the frozen official Qwen3-8B-Base router reaches
balanced accuracy $0.764\pm0.035$, macro-$F_1$ $0.775\pm0.032$, and routed MSE
$0.536$, and is the \emph{only} method that routes both oracle classes
correctly (trend recall $0.72$, periodic recall $0.81$).  Every alternative is
far below: a supervised MLP ($0.483$ balanced accuracy, MSE $1.147$), a
PatchTST encoder ($0.292$, MSE $1.971$), Chronos-T5-small/base ($0.133/0.264$,
MSE $2.06/1.04$), order-free marginal statistics ($0.475$, MSE $0.789$),
temporal hand features ($0.436$, MSE $1.129$), a random control ($0.022$, MSE
$2.885$), and constant-trend ($0.500$, MSE $0.590$).  Paired window-level
intervals show the Qwen router's MSE is significantly lower than every
baseline on most seeds (all three seeds for features, PatchTST, both Chronos
models, and random; two of three seeds for marginal, MLP, and constant-trend).
The transfer therefore comes from the language-pretrained representation, not
from being a better supervised classifier nor from the tested
time-series-pretrained representations.

\begin{table}[t]
\centering
\small
\begin{tabular}{lcc}
\toprule
Router & Raw oracle-hit & MSE \\
\midrule
Oracle (upper bound) & --- & $0.447$ \\
Constant trend & $\sim0.87$ & $0.500$ \\
Hand features & $0.133$ & $1.222$ \\
Random Qwen probe & $0.278$ & $1.098$ \\
\textbf{Frozen Qwen3-8B-Base} & $0.864$ & $0.538$ \\
\bottomrule
\end{tabular}
\caption{E4 imbalanced-set diagnostic (mixture/regime, 3 seeds pooled): raw
metrics are dominated by the majority (trend) oracle class and a constant-trend
policy already achieves high raw accuracy; see Table~\ref{tab:e4bal} for the
balanced test.}
\label{tab:e4}
\end{table}

\begin{table}[t]
\centering
\small
\begin{tabular}{lcccc}
\toprule
Representation & Balanced acc. & Macro-$F_1$ & MSE & Rec.\ T/P \\
\midrule
\textbf{Frozen Qwen3-8B-Base} & $\mathbf{0.764}$ & $\mathbf{0.775}$ & $\mathbf{0.536}$ & $0.72/0.81$ \\
Constant trend & $0.500$ & $0.333$ & $0.590$ & $1.00/0.00$ \\
Marginal statistics & $0.475$ & $0.355$ & $0.789$ & $0.93/0.02$ \\
MLP & $0.483$ & $0.549$ & $1.147$ & $0.64/0.32$ \\
Hand features & $0.436$ & $0.303$ & $1.129$ & $0.00/0.87$ \\
Chronos-T5-base & $0.264$ & $0.232$ & $1.038$ & $0.48/0.04$ \\
PatchTST & $0.292$ & $0.403$ & $1.971$ & $0.26/0.32$ \\
Chronos-T5-small & $0.133$ & $0.182$ & $2.056$ & $0.23/0.04$ \\
Random Qwen3-8B & $0.022$ & $0.037$ & $2.885$ & $0.04/0.00$ \\
\bottomrule
\end{tabular}
\caption{E4 balanced OOD test: trend/periodic-oracle balanced (60 each, 3
seeds pooled), official Qwen3-8B-Base.  Balanced accuracy, macro-$F_1$, routed
MSE, and per-class recall.  Qwen is the only method covering both oracle
classes.}
\label{tab:e4bal}
\end{table}

\paragraph{Imbalanced-set diagnostic (honest caveat).}
On the raw mixture/regime set (Table~\ref{tab:e4}), the oracle is trend-
dominated and raw metrics are dominated by the majority class: constant-trend
achieves $\sim$87\% raw oracle-hit and MSE $0.50$; the official Qwen router
reaches $86.4\%$ raw oracle-hit and MSE $0.538$ (better than random $0.278$ and
features $0.133$), and its paired MSE is not reliably below constant-trend
(significant only on one of three seeds).  We therefore do \emph{not} claim
that routing reduces MSE on this benchmark; the balanced test in
Table~\ref{tab:e4bal} is the evidence for compositional recognition transfer.

\paragraph{A pretrained \emph{time-series} encoder does not transfer OOD.}
To answer ``does this need a language model, or would a pretrained
\emph{time-series} foundation model do?'' we probe Chronos-T5
\citep{ansari2024chronos} with the identical linear-probe protocol.  Chronos
recognizes all five dynamics essentially perfectly in-domain ($99.7$--$100\%$),
yet on the balanced OOD test it routes at $0.13$--$0.26$ balanced accuracy and
MSE $1.0$--$2.1$, far below the frozen language model.  A representation
optimized for forecasting is therefore not what supports \emph{cross-dynamics}
recognition transfer; that capability comes specifically from language
pretraining.

\subsection{Real-world zero-shot routing (R6)}
We apply the router, trained only on synthetic clean dynamics, to 15 real-world
datasets (Table~\ref{tab:real}).  The frozen LLM router is best or near-best on
most datasets and is especially strong where feature routers collapse (e.g.,
Exchange rate: LLM MSE $0.037$ vs.\ feature $1.591$).  We also report honest
failures: on M4 macro series and some energy data the small $0.6$B router does
not beat the degenerate ``always-local'' baseline, showing that recognition
transfer requires sufficient scale and that routing helps most when dynamics
are heterogeneous.

\begin{table}[t]
\centering
\small
\begin{tabular}{lcccc}
\toprule
Dataset & Oracle & Best single & Feature & Frozen LLM \\
\midrule
ETTm1 & $0.011$ & periodic/$0.015$ & $0.521$ & $\mathbf{0.102}$ \\
ETTh1 & $0.028$ & periodic/$0.045$ & $0.581$ & $\mathbf{0.093}$ \\
Exchange & $0.017$ & periodic/$0.027$ & $1.591$ & $\mathbf{0.037}$ \\
Weather & $0.021$ & periodic/$0.025$ & $0.482$ & $\mathbf{0.319}$ \\
Traffic & $0.680$ & periodic/$0.883$ & $0.771$ & $0.913$ \\
Illness & $0.554$ & trend/$1.067$ & $4.620$ & $\mathbf{1.050}$ \\
\midrule
Electricity & $0.407$ & local/$0.855$ & $\mathbf{0.740}$ & $0.819$ \\
ETTh2 & $0.114$ & periodic/$0.211$ & $\mathbf{0.610}$ & $0.436$ \\
M4 (6 freqs) & --- & --- & best & \emph{not better} \\
\bottomrule
\end{tabular}
\caption{R6: zero-shot routing MSE on real datasets, 0.6B.}
\label{tab:real}
\end{table}

\paragraph{Scale improves real-data routing.}
The $8$B router (full table in Appendix~\ref{app:real8}) is stronger on the
same datasets: ETTh1 MSE $0.045$ (vs.\ feature $0.581$, $13\times$), ETTh2
$0.220$ (vs.\ $0.610$), ETTm1 $0.016$ (vs.\ $0.521$, $33\times$), Exchange
$0.028$ (vs.\ $1.591$, $58\times$), Weather $0.041$ (vs.\ $0.482$), Illness
$1.455$ (vs.\ $4.620$).  Across three router-training seeds on the same
evaluation windows, the LM router shows \emph{seven nominally significant}
gains (ETT$\times$4, Exchange, Weather, Illness), seven nominally significant
losses (Electricity, Traffic, and five M4 frequencies), and one inconclusive
comparison (M4-Hourly); no multiple-comparison correction is applied and these
real-data results are treated as an exploratory boundary check.  On the
periodic-dominated datasets (ETTm2, Weather, Exchange) the ``always-periodic''
degenerate baseline is near-optimal and the LLM router matches it without
harm; on Electricity and Traffic the feature router remains better; and on M4
both routers fail, showing that recognition transfer requires diverse
structure and does not hold universally.  The cross-family pattern is the
same: a GPT-2 router trained only on synthetic clean dynamics beats the
feature router on 8 of 15 real datasets, confirming that recognition transfer
is a property of language pretraining rather than of a specific model family.

\paragraph{The routing advantage holds at longer horizons.}
Table~\ref{tab:horizon} re-evaluates the $8$B router at horizons 24 and 48
(hidden states depend only on the context, so the same representations are
reused).  The frozen-LLM router beats the feature router on the same four of
five datasets at both horizons, tracking the best-single expert closely (e.g.,
Exchange H$=48$: $0.060$ vs.\ best $0.060$; ETTh1 H$=48$: $0.083$ vs.\ $0.071$),
and only Electricity remains favorable to features---exactly as at $H{=}16$.

\begin{table}[t]
\centering
\small
\begin{tabular}{llcccc}
\toprule
Dataset & Horizon & Oracle & Best single & Feature & Frozen LLM \\
\midrule
\multirow{2}{*}{ETTm1} & $24$ & $0.015$ & $0.019$ & $0.596$ & $\mathbf{0.028}$ \\
 & $48$ & $0.027$ & $0.034$ & $0.757$ & $\mathbf{0.051}$ \\
\multirow{2}{*}{ETTh1} & $24$ & $0.035$ & $0.056$ & $0.672$ & $\mathbf{0.061}$ \\
 & $48$ & $0.050$ & $0.071$ & $0.842$ & $\mathbf{0.083}$ \\
\multirow{2}{*}{Weather} & $24$ & $0.035$ & $0.044$ & $0.521$ & $\mathbf{0.123}$ \\
 & $48$ & $0.112$ & $0.148$ & $0.579$ & $\mathbf{0.317}$ \\
\multirow{2}{*}{Exchange} & $24$ & $0.026$ & $0.036$ & $1.944$ & $\mathbf{0.036}$ \\
 & $48$ & $0.046$ & $0.060$ & $2.497$ & $\mathbf{0.060}$ \\
\multirow{2}{*}{Electricity} & $24$ & $0.465$ & $0.792$ & $\mathbf{0.668}$ & $0.818$ \\
 & $48$ & $0.576$ & $0.910$ & $\mathbf{0.792}$ & $0.943$ \\
\bottomrule
\end{tabular}
\caption{R6: zero-shot routing MSE at horizons 24/48 (8B).  The LLM router
wins on 4/5 datasets at both horizons.}
\label{tab:horizon}
\end{table}

\subsection{Patch-level dynamic routing (E6)}
Routing per patch rather than per window lets the model switch experts as
dynamics change (Table~\ref{tab:e6}).  On regime-switching windows dynamic
routing reduces MSE by $33\%$ relative to the best static router and by $20\%$
on local windows, while remaining neutral on stationary windows---a clean
statement of when dynamic routing helps.

\begin{table}[t]
\centering
\small
\begin{tabular}{lccc}
\toprule
Window kind & Static best & Dynamic router & $\Delta$ \\
\midrule
Regime & $0.82$ & $0.55$ & $-33\%$ \\
Local & $1.00$ & $0.80$ & $-20\%$ \\
Stationary & $1.00$ & $1.00$ & $0\%$ \\
\bottomrule
\end{tabular}
\caption{E6: patch-level dynamic routing (normalized MSE).}
\label{tab:e6}
\end{table}

\section{Analysis: Why Recognition Transfers but Generation Fails}\label{sec:analysis}

We offer a functional decomposition.  Recognizing temporal structure requires
the model to map the observed context to a low-dimensional set of ``kinds''
(\emph{trend-like}, \emph{periodic}, \emph{noise-dominated}).  Language
pretraining produces representations in which such distinctions are linearly
separable---a reusable prior that generalizes to OOD dynamics because the
categories are compositional and the probe is simple.  In contrast, generating
future numeric values requires (a) a correct serialization/tokenization of
numbers, (b) precise extrapolation under distribution shift, and (c) exact
regression over a continuous target---precisely the skills language modeling
does not optimize.  The generation failures we observe (low parse coverage,
multiplicative MSE) are consistent with this decomposition, and the persistence
of the gap under constrained decoding isolates the regression/continuation
failure from formatting.

\paragraph{Formalization sketch (P2).}  Let $R_\theta$ be the recognition map
and $G_\theta$ the generation map of an LLM.  Treat classification as a
margin problem over $\{0,1\}^K$ and generation as iterated autoregression with
error accumulation rate $\rho$.  Recognition robustness requires only a margin
$\gamma>0$ at the linear probe, which pretraining supplies; generation error
grows as $\rho^H$ in horizon $H$.  The asymmetry then follows whenever
$\rho^H$ dominates the classification margin, which is exactly what we observe
as $H$ increases (horizon sweep in appendix).  This also predicts \emph{when}
routing helps: when dynamics are heterogeneous (mixtures, regime switches), the
oracle-expert variance is large and recognition-based routing captures most of
the gains; when dynamics are homogeneous, the best single expert is nearly
optimal and routing is neutral---consistent with E6.

\section{Limitations}\label{sec:limits}

Our controlled setup uses synthetic dynamics and three univariate experts; the
real-data routing inherits this simplicity.  The $8$B model is a custom
Wave-MoE checkpoint whose language trunk matches Qwen3, but which is not
identical to the vanilla $0.6$B/$1.7$B checkpoints, so the scale comparison is
approximate.  The M4 result shows that recognition transfer does not hold
universally at small scale.  Finally, we study frozen recognition; how
recognition evolves during fine-tuning is analyzed in the appendix but deserves
further study.

\section{Conclusion}\label{sec:conclusion}

The debate over whether language pretraining helps time-series forecasting is,
we argue, partly a debate over which capability---recognition or
generation---is being measured.  Under a controlled protocol, language
pretraining clearly helps forecasting, but the source of the benefit is a
reusable temporal-structure prior that transfers out-of-distribution through
recognition, while frozen numeric generation remains fragile.  This reframing
explains seemingly contradictory results, motivates using LLMs as temporal
expert routers rather than numeric forecasters, and opens a clean axis for
future analysis.

\bibliographystyle{iclr2027_conference}
\bibliography{refs}

\appendix
\section{Additional analyses}\label{app:extra}

\subsection{Pre- vs.\ post-fine-tuning representation}
\label{app:prepost}
Loading a LoRA adapter from E1 (fine-tuned on ETTm1 forecasting) onto the
frozen $8$B backbone and re-probing the \emph{same} synthetic windows shows
that forecasting fine-tuning does \emph{not} erode recognition: probe accuracy
rises from $0.967$ (frozen) to $0.983$ (post-FT), with the largest gains on
local ($0.95{\to}0.98$) and mixture ($0.92{\to}0.98$) windows.  The
representation itself shifts substantially (mean abs delta $0.69$; linear CKA
$0.45$; probe-direction cosine $0.08$), i.e., fine-tuning rewrites the
recognition geometry while preserving---even sharpening---the underlying
separability.

\subsection{Difficulty sweeps (context, noise, horizon)}
\label{app:asym}
See Table~\ref{tab:asym} in the main text.  Recognition degrades only mildly
under noise ($0.98$ at $\sigma{=}0.05$ to $0.91$ at $\sigma{=}0.2$) and
improves with horizon, while the generation/oracle MSE ratio stays $6$--$9\times$
and grows with horizon ($6.5\times$ at $H{=}8$ to $8.4\times$ at $H{=}32$).

\subsection{Tokenizer / template ablations}
\label{app:tok}
Prior work in this project (summarized in
Appendix~\ref{app:random_bpe}) showed that when a random-BPE tokenizer is
applied to numeric text, the model collapses to degenerate outputs, whereas
fixed-width two-decimal numeric tokens (our protocol) keep generation
parseable; the recognition probe is insensitive to this choice.  This isolates
the tokenization contribution to the generation failure channel.

\subsection{Full real-data table (8B)}
\label{app:real8}
See Table~\ref{tab:real8}.

\begin{table}[h]
\centering
\small
\begin{tabular}{lccc}
\toprule
Dataset & Feature MSE & Frozen-LLM MSE & Best single \\
\midrule
Electricity & $\mathbf{0.740}$ & $0.927$ & local/$0.855$ \\
ETTh1 & $0.581$ & $\mathbf{0.045}$ & periodic/$0.045$ \\
ETTh2 & $0.610$ & $\mathbf{0.220}$ & periodic/$0.211$ \\
ETTm1 & $0.521$ & $\mathbf{0.016}$ & periodic/$0.015$ \\
ETTm2 & $0.327$ & $0.114$ & periodic/$0.064$ \\
Exchange & $1.591$ & $\mathbf{0.028}$ & periodic/$0.027$ \\
Illness & $4.620$ & $\mathbf{1.455}$ & trend/$1.067$ \\
Traffic & $\mathbf{0.771}$ & $0.916$ & periodic/$0.883$ \\
Weather & $0.482$ & $\mathbf{0.041}$ & periodic/$0.025$ \\
M4-Daily & $\mathbf{1.244}$ & $1.826$ & local/$1.213$ \\
M4-Hourly & $214.8$ & $\mathbf{214.0}$ & periodic/$214.0$ \\
M4-Weekly & $\mathbf{1.419}$ & $2.524$ & local/$1.419$ \\
M4-Monthly & $\mathbf{1.218}$ & $2.140$ & local/$1.218$ \\
M4-Quarterly & $\mathbf{1.233}$ & $1.996$ & local/$1.233$ \\
M4-Yearly & $\mathbf{1.164}$ & $1.938$ & local/$1.164$ \\
\bottomrule
\end{tabular}
\caption{R6 full: official Qwen3-8B-Base zero-shot routing MSE across 15
datasets (seed 7, 300 windows per dataset).}
\label{tab:real8}
\end{table}

\paragraph{Real-data window sampling.}
For each non-M4 dataset we sample up to $300$ test windows by drawing
non-overlapping start positions uniformly at random without replacement
(context length $64$, horizon $16$), so windows within a dataset do not
overlap; each window is normalized using only its own context statistics.
ETT datasets use the standard $12/4/4$ month split; the remaining datasets use
a $70/30$ chronological split.  Because consecutive starts are drawn without
replacement from the available range, the effective stride is large relative
to the window length, and we report window-level bootstrap intervals
accordingly.

\paragraph{M4 evaluation protocol.}
M4 windows are sampled \emph{within the official training series} rather than
the official test series for two reasons: (i) the official competition
horizons ($6$--$48$, depending on frequency) differ from the uniform $H{=}16$
protocol used here, and (ii) using the short official test series would yield
far too few $C{=}64,H{=}16$ rolling windows for a stable evaluation.  This is
therefore a boundary-transfer diagnostic over many real macro series, not a
leaderboard evaluation.

\paragraph{Numerical parser.}
Generated completions are parsed with the regular expression
\texttt{(?<![0-9])[+-]?$\backslash$d+$\backslash$.$\backslash$d\{2\}} (any number
of integer digits, exactly two decimals, with a digit boundary), taking the
first $H$ matches.  An audit of all raw completions across models and seeds
found \emph{zero} generated values with magnitude $\geq 10$ and zero windows
where a one-digit-integer parser would have differed, so the reported parse
rates and MSEs are unaffected by the integer-width choice.

\subsection{Random-BPE tokenizer ablation}
\label{app:random_bpe}
We trained a BPE tokenizer on shuffled/random text and applied it to the same
numeric forecasting protocol.  Under this tokenizer, frozen generation
collapses (parse rate near zero), while the recognition probe is unaffected.
This confirms that the generation failure channel includes a tokenization
component---absent in our main protocol---and that recognition is robust to it.

\subsection{Statistical details}
\label{app:bootstrap}
We report 95\% nonparametric percentile bootstrap confidence intervals from
2,000 resamples.  In contrast to a seed-level analysis, the resampling unit is
the \emph{individual evaluation window}: all per-window outcomes are pooled
across seeds (3 seeds for Qwen3-8B; 5 seeds for Qwen3-0.6B) and resampled
with replacement.  The generation ratio is bootstrapped as a ratio of sums
($\sum\text{MSE}/\sum\text{oracle}$), matching the point estimate reported in
the main text.

For Qwen3-8B (pooled over 360 routing windows), pretrained routing accuracy is
$0.681$ [0.633, 0.728], random routing $0.286$ [0.239, 0.336], and feature
routing $0.133$ [0.100, 0.172]; the three intervals do not overlap.
Recognition is $0.963$ [0.951, 0.974] (pretrained) vs.\ $0.906$ [0.886, 0.923]
(random).  Routed MSE is $0.745$ [0.677, 0.815] (pretrained), $1.053$
[0.979, 1.129] (random), $1.222$ [1.148, 1.295] (feature), with an oracle of
$0.447$ [0.423, 0.471].  The window-level generation/oracle ratio is $5.28\times$
[3.22, 9.83].

For Qwen3-0.6B (pooled over 600 routing windows), pretrained routing is
$0.337$ [0.298, 0.375] vs.\ random $0.217$ [0.183, 0.250] vs.\ feature $0.162$
[0.132, 0.192]; recognition $0.984$ [0.977, 0.989] vs.\ $0.789$ [0.769, 0.809];
and the generation/oracle ratio is $9.64\times$ [7.52, 12.33].
The window-level intervals confirm the same ordering at both scales.

\end{document}
